\section{The Nachtenschein of Classic Liberalism}

Having come across a remark in AKC\footnote{\url{http://akc.satishankar.com/2012/09/selected-letters-of-ak-coomaraswamy.html}}`s letters that the really cultured and spiritual European does not have a peer in their Eastern counterparts, I returned to a volume of Wilhelm Humboldt\footnote{\url{https://en.wikipedia.org/wiki/Wilhelm\_von\_Humboldt}}`s collected letters and essays, excerpted by subject. Although one can tell that the writing was not in English (the German lumbers through the translations, probably because English doesn't have words for some of these ideas), and although he was not per se a traditionalist or even a religious man, it contains a testimony to the power of spiritual perception that was once much more common in Europe -- that of the civilized Christian, in this case, looking at the unfolding modern Age. In this excerpt, he mentions the \emph{Bhagavad-Gita}, which (incidentally) he was much interested in. In this selection, he defends the idea that poetry and philosophy jointly can express (at their highest) the same truths that religion has been accustomed to deal with (Matthew Arnold thought the same thing \& incidentally, John Ruskin thought that the ``pathetic fallacy'' was only a fallacy in the hands of lesser masters\footnote{\url{http://www.ourcivilisation.com/smartboard/shop/ruskinj/}}). My question would be, isn't religion just exactly that, what they describe? Why re-invent what is already there? Isn't religion the relationship of Life between Poetry, myth, and philosophy? What is astonishing in reading his letters is how close he comes to the truth by his own native powers of cultured perception, \& yet how far he is from the kind of mental clarity that AKC embodied, or that Traditional studies aims at. My conclusion is that he was indeed a cultured, \emph{but not a spiritual}, European, \& that this flaw has largely to do with the decline of Christianity within its traditional pale\footnote{\url{https://en.wikipedia.org/wiki/The\_Pale}}. One cannot help but think that the greatness of these ``classical liberals'' springs from their \emph{Nachtenshein} of Christianity, which still retained enough magic to cast a spell over the Iron Age which they were unwittingly hatching.

\begin{quotationx}
Poetry and philosophy grow out of the same soil; they come from the highest and deepest levels of mankind, and the difference between a genuine philosophical poem and one which does not deserve the name lies in the question whether the two sources are represented as truly joint or whether they have been mechanically superimposed upon one another.

It is a prerogative of poetry to claim man's total undivided nature and to lead him to the point where his finite nature loses itself in intuitions of infinity. Poetry deserves its name only when it achieves this. Therefore nothing is excluded from its proper field -- no object, no genre, not the plainest elegaic, not the most lightheartedly joyous nor the more irrepressibly or capriciously comic verse. For feeling itself, partly in its own movements, but especially when it is purified by the aesthetic sense whose activity is always stimulated in man at the first musical note, carries a relationship to infinity. An artform knows no barriers except those placed by its own definition. But the true secret lies in the creative imagination within which all art operates and which it controls. By its magic power in a mode very much like that described in the \emph{Bhagavad-Gita} it so destroys finite nature in its substance and so preserves it in its forms that, living in the midst of sense involvements, it dissolves all sentient emotion into pure, ideal-oriented perception, just as the dogma of renunciation and immersion dissolves the liveliest action into non-action. What Krishna says of creatures, namely that they meet by remain unknown to each other, like sudden apparitions (II, 29) is also and characteristically true of any genuine poem. It stands before us, and not footprints tell us whence it came. It requires certification from elsewhere, therefore, and the invocation to a higher power is a need felt as natural by every poet, insofar as he is not -- as is the author here under consideration --- imbued with the feeling that he is carrying such certification from on high within himself.

If poetry, then, is to combine with philosophic ideas in a worthy manner, the latter must be of the sort that they too could not have come about without an invisible capacity for inner inspiration. The fire and loftiness of poetry must look necessary for calling forth truth from the depths of the spirit; philosophic dogma must not seek a poetic garment like a borrowed ornament but must pour itself out in voluntary, free rhythms, moved by an inner urge; it must move within poetry as though in a natural, indigenous form. But this can only be the case if the philosophic ideas go back to the point where the rationalizing intellect has to give up its effort to develop effects from causes; the point where truth flames forth by the mere purification and orientation of spirit, by the removal of all dialectic semblances, from an intensification of pure self-awareness. In this realm, where the poet feels strong enough to maintain the nature of truth even in the midst of ardent poetic imagination --- imbued with the feeling that he is carrying such certification from on high here alone is found the true philosophical poem.
\end{quotationx}



\flrightit{Posted on 2012-10-12 by Logres}

\begin{center}* * *\end{center}

\begin{footnotesize}
\begin{sffamily}
\texttt{Cassiodorus on 2012-10-13 at 12:04 said:}

The nature of the relationship between Christianity and classical liberalism is a subject that has taxed me ever since I had my first mature thoughts. It would seem that from the perspective of Tradition, the metaphysical order of Christendom and ``Enlightenment'' rationalism are utterly incompatible- Pope St Pius X referring to modernity as a ``many headed hydra'', the heresy of heresies. On the other hand, many proponents of ``conservatism'' today seem to argue that classical liberalism represents the flowering of Christian principles, being natural partners in which liberty \emph{depends} on the Creator and a transcendent moral order.

The Traditionalist, Titus Burckhardt, commented in his paper ``What is Conservatism'', that conservatism was born the moment that Tradition came to an end in the West.

\hfill

\texttt{Logres on 2012-10-13 at 22:12 said:}

That's actually a very helpful observation -- thank you Cassiodorus. I think the sin of the classical liberals was to fail to see the source and origin of their celebrated liberty -- it ``looked'' to them like this was possible for all people, everywhere, rather than being a happy flourishing for certain climes and seasons. As an ideology, it is antithetical to Tradition -- and it pains me to come to that conclusion, because I was bred on such things.

\hfill

\texttt{Graham on 2012-10-14 at 01:01 said:}

``Isn't religion the relationship of Life between Poetry, myth, and philosophy?''

And this relationship is manifested in prayer. Take the litany of the blessed virgin, for example, or the canticle of creatures. These prayers flow from contemplation, and back towards it.

\hfill

\texttt{Logres on 2012-10-14 at 21:03 said:}

I suppose so. I recently read advice on prayer to focus on the chest, rather than the heart, and let the heart be attracted to the warmth of the prayer. High liturgical poetry has always drawn me, even when outside its circle. Repetition is now acquiring a peculiar, rare beauty for me.

\hfill

\texttt{Graham on 2012-10-14 at 21:33 said:}

``Why re-invent what is already there?''

I took a poetry class once, there was a poet who came to speak, he said ``poetry is the closest thing to prayer.'' I thought, well, why not just pray ?

He then recited a poem about getting drunk and losing control of his bowels.

\hfill

\texttt{Graham on 2012-10-14 at 21:40 said:}

Repetition: yes, litanies are beautiful. Time stops during a good litany.

\hfill

\texttt{Logres on 2012-10-16 at 06:19 said:}

Interesting. Considering the poetry of the Scriptures is first rate in the Wisdom books, I have a hard time understanding this also -- or pick some Latin hymns. On a deeper level, poetry and prayer are probably the same thing, with sacred poetry being actual prayer. I'm not sure anyone has really done this well in the West, although I haven't read John of the Cross' poetry... .yet.

\hfill

\texttt{Cologero on 2012-10-18 at 08:17 said:}

Jesus \& his Apostles \& Disciples were all Artists

A Poet a Painter a Musician an Architect : the Man

Or Woman who is not one of these is not a Christian

William Blake


\hfill

\texttt{Logres on 2012-10-19 at 23:10 said:}

Humboldt, I think, was arguing this in his excerpt I just put up -- I am thinking that he embodies the old ideal of cultured man quite well, with the exception that he had no external help from Traditional doctrine -- but he yearns for it. Unfortunately, brilliant though he was, he didn't write great poetry. But he tried it anyway... 

\hfill

\end{sffamily}
\end{footnotesize}

